
\def\PUDcodigo{ACE021}

\if\COMPILAR0
   \def\PUDcodacad{04507.59}
\else
   \def\PUDcodacad{04506.57}
\fi

\def\PUDdisciplina{Educação Física}

\def\PUDsemestre{Opt}

\def\PUDhoras{40}

%NOVOS ---------------------------------------------------
\def\PUDhorasTeoria{20}
\def\PUDhorasPratica{20}

\def\PUDinterECA{---}

\def\PUDatuaisECA{---}

\def\PUDrecursos{Práticas esportivas realizadas no ginásio, campo e/ou academia, com materias esportivos do campus. Eventualmente, podem ser realizadas aulas em sala de aula com projetor multimídia, quadro branco e pincel.}

%----------------------------------------------------------

\def\PUDcreditos{4}

\def\PUDprerequisito{---}

\def\PUDpreacad{---}

\def\PUDobjetivos{Oportunizar aos alunos o contato com os mais variados tipos de atividades esportivas, visando capacitá-los para uma atitude crítica, saudável e ética; Aprofundar o entendimento dos alunos a cerca da utilização do  corpo  no  processo  de  trabalho de uma maneira saudável, sem causar danos à saúde; Estimular a vivência do lazer, do trabalho e da saúde através de práticas esportivas; Capacitar os alunos a uma prática de atividade física autônoma e efetiva; Desenvolver as capacidades físico-motoras e psicológicas, visando o trabalho em equipe, almejando um objetivo comum.}

\def\PUDementa{Proporcionar atividades por meio de práticas com predominância às de natureza desportiva, preferencialmente as que conduzam à manutenção e aprimoramento da aptidão física, à conservação da saúde, à integração do estudante ao meio universitário, à consolidação de comunidade e nacionalidade.}

\def\PUDconteudo{ &  Introdução à disciplina. 
\\ 
 &  Prática e desenvolvimento de esportes.
\\ 
 &  LER e DORT no dia a dia dos alunos.
\\
 &  Prática de qualidade de vida e lazer. }

\def\PUDmetodo{Aulas práticas de atividades esportivas ao longo da disciplina.}

\def\PUDavalia{A avaliação será feita de forma contínua ao longo das atividades realizadas ao longo da disciplina.}

\def\PUDbibbasica{ &  \href{www.biblioteca.ifce.edu.br/index.asp?codigo_sophia=42081}{Couto}, Hudson de Araújo.  Gerenciando a LER e os DORT nos tempos atuais.  Belo Horizonte, MG: Ergo, 2007.  492p.  ISBN: 978859959028.
\\
&   \href{www.biblioteca.ifce.edu.br/index.asp?codigo_sophia=39440}{Vitte}, Claudete de Castro Silva; Keinert, Tânia Margarete Mezzomo.  Qualidade de vida, planejamento e gestão urbana: discussões teórico-metodológicas.  Rio de Janeiro, RJ: Bertrand Brasil, 2009.  310p.  ISBN: 9788528613865. 
\\ 
 &   \href{www.biblioteca.ifce.edu.br/index.asp?codigo_sophia=23862}{Salvador}, César Coll.  Desenvolvimento psicológico e educação: psicologia da educação escolar.  v.2.  2º ed.  Porto Alegre, RS: Artmed, 2007. 472p.  ISBN: 9788536302287. }

\def\PUDbibcomplementar{ &  \href{www.biblioteca.ifce.edu.br/index.asp?codigo_sophia=57561}{Finck}, Silvia Christina Madrid.  A Educação Física e o Esporte na Escola cotidiano saberes e formação.  E-book.  194p.  ISBN: 9788582120330.
\\
&  \href{www.biblioteca.ifce.edu.br/index.asp?codigo_sophia=59645}{Weineck}, Jurgen.  Anatomia aplicada ao esporte.  18ª ed.  E-book.  372p.  ISBN: 9788520432044.  	
\\ 
 &  \href{www.biblioteca.ifce.edu.br/index.asp?codigo_sophia=58745}{Puleo}, Joe.  Anatomia da Corrida: Guia Ilustrado de Força, Velocidade e Resistência para Corrida.  E-book.  202p.  ISBN: 9788520431627.
\\
 & \href{www.biblioteca.ifce.edu.br/index.asp?codigo_sophia=56563}{Evans}, Nick.  Anatomia da Musculação.  E-book.  204p.  ISBN: 9788520426258.
\\
 & \href{www.biblioteca.ifce.edu.br/index.asp?codigo_sophia=58744}{McLeod}, Ian.  Anatomia da Natação. E-book. 204p.  ISBN: 9788520431177. }

\def\PUDautor{Samuel Vieira Dias}

\def\PUDdata{2017-07-07}

\def\PUDrevisao{1}

\def\PUDrevisor{Samuel Vieira Dias}

\def\PUDdatarevisao{2017-07-11}

\PUDcorpo
